\subsubsection{Edmonds-Karp (max flow)}
\begin{lstlisting}[style=mystyle, language=C++]
// TC: O(N * M^2)
// usa: max flow simple; sirve como baseline o para grafos pequenos
#include <bits/stdc++.h>
using namespace std;

const long long INF = 1e18;

long long edmondsKarp(int n, vector<vector<pair<int,long long>>>& adj, int s, int t){
	long long flow = 0;
	while(true){
		vector<int> parent(n, -1);
		queue<int> q;
		parent[s] = s;
		q.push(s);
		while(!q.empty() && parent[t] == -1){
			int v = q.front(); q.pop();
			for(auto [u, cap] : adj[v]){
				if(cap > 0 && parent[u] == -1){
					parent[u] = v;
					q.push(u);
				}
			}
		}
		if(parent[t] == -1) break;
		long long add = INF;
		for(int v=t; v != s; v = parent[v]){
			for(auto [u, cap] : adj[parent[v]]){
				if(u == v) add = min(add, cap);
			}
		}
		for(int v=t; v != s; v = parent[v]){
			for(auto& [u, cap] : adj[parent[v]]){
				if(u == v) cap -= add;
			}
			for(auto& [u, cap] : adj[v]){
				if(u == parent[v]) cap += add;
			}
		}
		flow += add;
	}
	return flow;
}

int main(){
	// 4-node max flow graph
	vector<vector<pair<int,long long>>> adj(4);
	auto add = [&](int u, int v, long long c){
		adj[u].push_back({v,c});
		adj[v].push_back({u,0});
	};
	add(0,1,10); add(0,2,5); add(1,2,15); add(1,3,10); add(2,3,10);
	cout << "max flow = " << edmondsKarp(4, adj, 0, 3) << "\n";  // 15
	return 0;
}
\end{lstlisting}
